%%%%%%%%%%%%
%
% $Autor: Wings $
% $Datum: 2019-03-05 08:03:15Z $
% $Pfad: TemplateSensor $
% $Version: 4250 $
% !TeX spellcheck = en_GB/de_DE
% !TeX encoding = utf8
% !TeX root = filename 
% !TeX TXS-program:bibliography = txs:///biber
%
%%%%%%%%%%%%
\chapter{Beispielimplementierung: Temperatursensor DS28B20}
	
Die Implementierung der Klasse \PYTHON{CDS18B20Sensor} folgt der Struktur des Interfaces \PYTHON{ISensor} und definiert die drei Pflichtmethoden:
	
\begin{itemize}
  \item \PYTHON{init()}:
  
		Initialisiert den 1-Wire-Bus und überprüft, ob genau ein DS18B20-Sensor mit dem angegebenen Pin verbunden ist.  
		Wenn kein Sensor gefunden wird oder mehrere Geräte auf der gleichen Leitung erkannt werden, wird die Initialisierung abgebrochen und ein entsprechender Fehlercode zurückgegeben.  
		
		\textbf{Rückgabewerte:}
		
		\begin{itemize}
			\item \PYTHON{SensorError::SUCCESS} – genau ein Sensor wurde erkannt und erfolgreich initialisiert. 
			\item \PYTHON{SensorError::NOT\_INITIALIZED} – kein Sensor auf dem Bus gefunden. 
			\item \PYTHON{SensorError::READ\_ERROR} – mehrere Sensoren auf derselben Leitung erkannt (ungültige Konfiguration). 
		\end{itemize}
		
    \item \PYTHON{update()}:
    
		Führt eine neue Temperaturmessung durch.  
		Hierzu wird eine Konvertierung über die Bibliothek \PYTHON{DallasTemperature} ausgelöst und das Ergebnis ausgelesen. Die Methode prüft, ob der Sensor initialisiert wurde und ob der gelesene Wert plausibel ist.  
		Werte au{\ss}erhalb des gültigen Bereichs (\(-55\,^\circ C\) bis \(125\,^\circ C\)) oder ein \PYTHON{DEVICE\_DISCONNECTED\_C}-Signal werden als Fehler interpretiert.  
		
		\textbf{Rückgabewerte:}
		
		\begin{itemize}
			\item \PYTHON{SensorError::SUCCESS} – Temperatur erfolgreich gemessen und gespeichert. 
			\item \PYTHON{SensorError::NOT\_INITIALIZED} – der Sensor wurde vor dem Messvorgang nicht initialisiert. 
			\item \PYTHON{SensorError::READ\_ERROR} – ungültiger Messwert (NaN oder au{\ss}erhalb des gültigen Bereichs). 
			\item \PYTHON{SensorError::SENSOR\_DISCONNECTED} – der Sensor ist physisch getrennt oder liefert kein Signal. 
		\end{itemize}
		
		\item \PYTHON{getValue()}:
		
		Gibt den zuletzt gespeicherten Temperaturwert über eine Referenzvariable zurück, ohne eine neue Messung durchzuführen.  Diese Methode dient ausschlie{\ss}lich dem Abruf des letzten gültigen Messwerts.  
		
		\textbf{Rückgabewerte:}
		\begin{itemize}
			\item \PYTHON{SensorError::SUCCESS} – ein gültiger Messwert ist verfügbar und wurde erfolgreich übergeben. 
			\item \PYTHON{SensorError::NOT\_INITIALIZED} – der Sensor wurde noch nicht initialisiert, daher steht kein Messwert zur Verfügung. 
		\end{itemize}
	\end{itemize}
	
{
  \captionof{code}{Klasse \PYTHON{CDS18B20Sensor}: Implementierung des Interfaces \PYTHON{ISensor} für Temperaturmessung}
  \ArduinoExternal{}{../../Code/MKRWAN1310/SoftwareArchitecture/CDS18B20Sensor.h}
}
	
{

  \captionof{code}{Implementierung der Methoden des Temperatursensors \PYTHON{CDS18B20Sensor}}
  \ArduinoExternal{}{../../Code/MKRWAN1310/SoftwareArchitecture/CDS18B20Sensor.cpp}
}
	